\documentclass[10pt,a4paper]{amsart}
\usepackage[margin=19mm]{geometry}
\usepackage{amsmath,amssymb,mathtools}
\usepackage[scr=rsfs,cal=euler]{mathalfa}
\usepackage{xfrac}
\usepackage[unicode,hidelinks,bookmarks=false]{hyperref}
\newcommand{\ie}{i.e.\ }
\newcommand{\cf}{cf.\ }
\newcommand{\resp}{respectively\ }
\newcommand{\Subsection}{\subsection}
\providecommand{\nobreakdash}{\nobreak\hbox{-}}
\setlength{\emergencystretch}{2em}
\multlinegap0pt
\usepackage{amssymb}
\usepackage{tikz}
\usepackage[shortlabels]{enumitem}
\newtheorem{lem}{Lemma}
\newtheorem{theorem}{Theorem}
\newcommand{\C}{\mathbb C}
\newcommand{\cI}{\mathcal I}
\newcommand{\g}{\mathfrak}
\newcommand{\sseq}{\subseteq}
\title{Radial restriction of spherical functions on supergroups}
\author{Mitra Mansouri, Hadi Salmasian}
\date{Source: arXiv:2504.19102v1 (CC BY 4.0)}
\begin{document}
\maketitle

\noindent\emph{Source and licence.} Radial restriction of spherical functions on supergroups. Version arXiv:2504.19102v1 (CC BY 4.0), distributed by arXiv under the Creative Commons Attribution 4.0 International licence, http://creativecommons.org/licenses/by/4.0/, as recorded in the arXiv licence metadata for this exact version and shown on the abstract page https://arxiv.org/abs/2504.19102v1. Full licence notice: \texttt{licenses/arxiv-2504.19102-CC-BY-4.0.txt}.\par\smallskip

\noindent\textit{Editorial scope.} Counted unit: the article's theorem I (source0.tex lines 1147--1154). The article's complete original proof of it is delivered with it (source0.tex lines 1156--1172, one \texttt{proof} environment of 17 lines). No mathematics is written, repaired or re-derived: the statement and the proof are the article's own lines, in the article's own notation, and no retained line is substituted or re-flowed.  Retained uncounted as the source-local context the proof needs: lines 1047--1064; lines 1081--1098.  Nothing was omitted from any retained span, so the package carries no omission record.  No retained line contains a citation, so the wrapper carries no bibliography and no bibliography entry.  No retained line contains a figure environment, an image-inclusion command or drawing code, so no image is delivered and the package declares no asset dependency.  Editorial formatting only, in addition: an AMS \texttt{amsart} preamble that restates the article's own declarations and macros used by the retained lines, namely the environments \texttt{lem}, \texttt{theorem} on separate counters, together with the standard packages those lines need.  The retained environments print in this document as Lemma 1, Theorem 1 in order of appearance, whereas the article prints the same items with the numbering of its own full document.  The complete native source of the article as distributed in the arXiv e-print for this exact version is bundled unchanged as \texttt{sources/177145/source0.tex}.
\begin{lem}\label{firstpart}
For every 
\[k_0=\begin{bmatrix}
0&0&0\\
0&\alpha&\beta\\
0&\gamma&-\alpha
\end{bmatrix}
\in \g k_{\bar{0}},
\] the following relations hold in $U(\g g)$:
\begin{itemize}
\item[\rm (i)] 
$p^nk_0=k_0p^n$.
\item[\rm (ii)]
$p^nek_0=k_0p^ne+p^n(\alpha e+\beta f)$.
\item[\rm (iii)] $p^nfk_0=k_0p^nf+p^n(\gamma e-\alpha f)$.
\item[\rm (iv)] $p^nefk_0=k_0p^nef+\beta k_1 p^n-\gamma k_2 p^n$.
\end{itemize}
\end{lem}

\begin{lem}\label{Second part}
Let $v=(a,b)\in\C^2$. Then the following relations hold in $U(\g g)$:
\begin{itemize}
\item[\rm (i)]
$p^nv_{\mathfrak k}=v_{\mathfrak k}\alpha_{n}(p)+\beta_{n-1}(p)v_{\mathfrak p}$.
\item[\rm (ii)]
$p^nev_{\mathfrak k}=-v_{\mathfrak k}\alpha_{n}(p)e+ak_2\beta_{n-1}(p)-bk\beta_{n-1}(p)+b\beta_{n-1}(p)ef-bp^{n+1}$.

\item[\rm (iii)]
$ p^nfv_{\mathfrak k}=-v_{\mathfrak k}\alpha_{n}(p)f-a\beta_{n-1}(p)ef-bk_1\beta_{n-1}(p)+ap^{n+1}$.
\item[\rm (iv)]
$
p^nefv_{\mathfrak k}=D-bk_1\beta_{n-1}(p)e-b\beta_{n-1}(p)f+p^{n+1}v_{\mathfrak p}-ae'\alpha_{n}(p)-\beta_{n-1}(p)e$, 
where \[
D:=v_{\mathfrak k}\alpha_{n}(p)ef-ak_2\beta_{n-1}(p)+bk\beta_{n-1}(p)f.
\]
\end{itemize}
\end{lem}

\begin{theorem}\label{basis of I}
The following vectors in $U(\g g)$ form  a basis of $\cI$:

\begin{itemize}
\item[\rm (i)] $z^mp^ne$ and $z^mp^nf$ where $m,n\geqslant 0$,
\item[\rm (ii)] $z^m\beta_{n-1}(p)ef-z^mp^{n+1}$ where $m,n\geqslant 0$, with the convention $\beta_{-1}(x):=0$,\item[\rm (iii)] $z^mk^{r_1}k_1^{r_2}k_2^{r_3}(e')^{r_4}(f')^{r_5}p^{s_1}e^{s_2}f^{s_3}$ where $m,r_1,r_2,r_3,s_1\geq 0$, $r_4,r_5,s_2,s_3\in\{0,1\}$, and 
 at least one of $r_1,\cdots , r_5\geq 1$.
\end{itemize}\end{theorem}

\begin{proof}
We remark that $z$ is in the centre of $U(\g g)$.
Let $\mathcal B_\mathrm{(i)},\mathcal B_\mathrm{(ii)}$
and 
$\mathcal B_\mathrm{(iii)}$ denote the sets of vectors defined in (i)-(iii) above, respectively. Then from Lemma~\ref{firstpart}(ii) it follows that $\mathcal B_\mathrm{(i)}\sseq \cI$. From Lemma~\ref{Second part}(iii) it follows that $\mathcal B_\mathrm{(ii)}\sseq \cI$. It is also clear that $\mathcal B_\mathrm{(iii)}\sseq\cI$.  

To prove that $\mathcal B_\mathrm{(i)}\cup\mathcal B_\mathrm{(ii)}\cup\mathcal B_\mathrm{(iii)}$ spans $\cI$, note that elements of $U(\g g)\g k$ that are not in $\g kU(\g g)$ are linear combinations of the
 monomials in the PBW basis that, modulo a factor of $z^m$,  appear on  the left hand sides of the relations in Lemmas~\ref{firstpart} and~\ref{Second part}. 
From these lemmas it follows that the latter monomials are in the span of $\mathcal B_\mathrm{(i)}\cup\mathcal B_\mathrm{(ii)}\cup\mathcal B_\mathrm{(iii)}$.

Finally, we prove that the vectors
in $\mathcal B_\mathrm{(i)}\cup\mathcal B_\mathrm{(ii)}\cup\mathcal B_\mathrm{(iii)}$
 are linearly independent. Consider any  linear dependence relation between the elements of 
$\mathcal B_\mathrm{(i)}\cup\mathcal B_\mathrm{(ii)}\cup\mathcal B_\mathrm{(iii)}$.  
By organizing the dependence relation according to the powers of $z$, and factoring out the common power $z^m$, we can assume that in all of the monomials $m=0$. 
  Let $c_n$ denote the coefficient of $\beta_{n-1}(p)ef-p^{n+1}$ for $n\geq 0$, and set $N:=\max\{n:c_n\neq 0\}$. Recall that $\deg(\beta_{N-1})=N-1$. Now with the PBW basis of $U(\g g)$ in mind, note that the summand $c_Np^{N-1}ef$ cannot get canceled in the linear dependence relation. This is a contradiction. 
\end{proof}

\end{document}
